\section{Related Work}
\label{sec:related}

\para{MLPs as key-value memories.}
\citet{geva2021transformer} establish that transformer feed-forward layers operate as key-value memories, where each neuron's key vector detects a specific input pattern and its value vector stores an associated output distribution.
\citet{geva2022transformer} extend this view by showing that individual MLP sub-updates (the product of a neuron's activation and its value vector) promote interpretable token concepts when projected to vocabulary space.
\citet{dai2022knowledge} identify ``knowledge neurons'' whose activation directly correlates with specific factual knowledge.
Our work builds on this foundation but asks a different question: rather than examining individual neurons in isolation, we test whether the \emph{aggregate} MLP input-to-output mapping can be approximated by a shallow linear program in token space.

\para{Factual recall and model editing.}
\citet{meng2022locating} use causal tracing to localize factual associations to specific MLP layers and develop rank-one model editing (ROME) by modifying MLP value matrices.
\citet{merullo2023language} show that MLPs implement simple vector arithmetic for relational tasks---the MLP produces a content-independent update that can be patched across examples.
These results suggest that at least some MLP computations are simple token-direction programs.
We go further by measuring how much of the \emph{total} MLP computation (not just factual recall) this description captures, and we find that the picture is more nuanced: the simple program works well in late layers but breaks down in middle layers where the most complex computation occurs.

\para{Linear representations.}
The linear representation hypothesis~\citep{park2024linear} holds that concepts in language models are encoded as directions in activation space.
\citet{marks2023geometry} provide evidence that truth values are linearly represented.
\citet{belrose2023eliciting} improve on the logit lens~\citep{nostalgebraist2020logit} by learning per-layer affine transformations for better intermediate decoding.
Our use of the unembedding matrix to project MLP activations into vocabulary space follows this line of work.
Unlike prior studies that focus on whether specific concepts are linearly decodable, we ask whether the MLP's computation itself can be linearly approximated \emph{in token space}.

\para{Superposition and sparse autoencoders.}
\citet{elhage2022toy} demonstrate that neural networks represent more features than they have dimensions through superposition.
\citet{cunningham2023sparse} train sparse autoencoders (SAEs) on MLP activations to discover interpretable features that go beyond individual neurons.
These results suggest that the token-direction basis may be insufficient to describe MLP computation because features are superposed across directions.
Our random-direction control experiment provides direct evidence for this: the vocabulary basis has no special advantage over random bases, consistent with MLP features being distributed across many directions rather than aligned with token embeddings.
